This section describes the database and web server configuration
procedure. It has been realised on a {\em Debian 13} system and is
intended for system administrators. It assumes that you have already
installed the webamc package, following the instructions in
Section~\ref{sec:install}.
%%%%%%%%%%%%%%%%%%%%
\subsection{Database creation}
Once package has been installed, a database must be created. Right now
the database cannot be created with the webamc package, so we will
interact with the server using the
\name{psql} client tool. First connect to the database server using
\name{postgres} identity.
\begin{bash}
$ sudo -u postgres psql
\end{bash}%$

In the postgresql interpreter, create the database, a database user
and grant this user all privileges on the created database:
\begin{sql}
create user webamc_user password 'webamc_password';
create database webamc;
alter database webamc owner to webamc_user;
grant all privileges on database webamc to webamc_user;
\end{sql}
%%%%%%%%%%%%%%%%%%%%
\subsection{Database initialisation}
The next step is to initialise the database: we have to create tables
and put some initial data in these. This is done via the \name{webamc}
tool provided by the package.

The configuration of the webamc command-line tool and web interface is
done via a JSON file defining several configuration parameters (e.g.,
database access parameters). An example file can be found in the
\file{examples/cfg} directory of the
repository. Section~\ref{sec:config} describes the different
configuration parameters.

By default, the configuration file is loaded from the
\file{.local/share/webamc/config.json} file in the home directory of
the user.
%%%%%%%%%%
\subsubsection{Tables creation}
First we need to initialise the database structure and create all its
table:
\begin{bash}
$ webamc --config-file /path/to/my/config.json init
\end{bash}%$
The command normally silently terminates. An error will likely be due
to a misconfiguration of the database access.
%%%%%%%%%%
\subsubsection{CSV files loading}
The \command{webamc loadcsv} command can be used to load records in
the database. To use it, simply type this command followed by the path
of one (or more) CSV file(s).
\begin{bash}
$ webamc --config-file /path/to/my/config.json loadcsv *.csv
\end{bash}

The \command{webamc genexamples} command creates a bunch of example
CSV files in the directory passed as argument.
\begin{bash}
$ mkdir csv && webamc genexamples csv
$ ls csv
00-usr.csv  01-cas_auth.csv  02-basic_auth.csv  03-tag.csv
04-grp.csv  05-usr_grp.csv   06-tbl.csv         07-admin.csv
$ webamc --config-file /path/to/my/config.json loadcsv csv/*.csv
\end{bash}

For a basic setup, you will at least need to insert a user with
administration privileges and an authentication mode so that she/he
can access the web interface. CSV files can also be loaded from the
web interface. Section \ref{sec:csv} describes the structure of CSV
files.
%%%%%%%%%%%%%%%%%%%%
\subsection{Web server execution and configuration}
Once the webamc package is installed, you can simply launch webamc's
web server via \texttt{uvicorn} and bind it to TCP port, e.g., 8008.
\begin{bash}
$ uvicorn webamc.www.app:app --port 8008
\end{bash}%$

You can also change set the {\tt CONFIG\_FILE} environment variable to
specify the configuration file to load:

\begin{bash}
$ CONFIG_FILE=/path/to/my/config.json uvicorn webamc.www.app:app --port 8008
\end{bash}%$

The \file{webamc.service} file in the root directory of the repository
contains a minimal service specification to interact with the service
via systemd.
%%%%%%%%%%%%%%%%%%%%
\subsection{Apache configuration}
Now that our server listens on the 8008 port, we configure apache to
redirect requests addressed to our server to this port.  The first
step is to enable the \texttt{proxy\_http} module of apache:

\begin{bash}
$ sudo a2enmod proxy
$ sudo a2enmod proxy_http
\end{bash}%$$

Then we have to add one line in the configuration of apache's virtual
host to pass our server all requests on paths starting by, e.g.,
\file{/webamc}.

\begin{bash}
ProxyPass /webamc http://localhost:8008/
\end{bash}

Then restart apache:
\begin{bash}
$ sudo systemctl restart apache2.service
\end{bash}
